\documentclass[11pt]{article}
\usepackage[margin=1in]{geometry}
\usepackage{amsmath,amssymb,amsthm}
\usepackage{enumitem}
\usepackage{graphicx}
\usepackage[hidelinks]{hyperref}

\title{Full Unimodular Cyclicity via a Bilateral Fixed-Space Dilation}
\author{Candidate full solution for the FA/Banach discovery run}
\date{23 July 2026}

\newtheorem{theorem}{Theorem}
\newtheorem{lemma}[theorem]{Lemma}
\newtheorem{corollary}[theorem]{Corollary}
\newcommand{\Fix}{\operatorname{Fix}}
\newcommand{\cL}{\mathcal L}
\newcommand{\Torus}{\mathbb T}

\begin{document}
\maketitle

\section*{Source problem and answer}

The source is Jochen Gl\"uck, \emph{A note on the fixed space of positive
contractions}, arXiv:2108.05174 \cite{GlueckFixed}.  Its Open Problem~3.4 on
page~6 asks whether Corollary~3.3 remains true when the eigenvalue
\(\lambda\) is an arbitrary unimodular number rather than a root of unity.
Thus, for a positive contraction \(T\) on a monotonically complete complex
Banach lattice with the Fatou property, the question asks whether
\[
 \dim\ker(\lambda-T)\leq \dim\ker(\lambda^k-T)
 \qquad(k\in\mathbb Z)
\]
whenever \(|\lambda|=1\).

The answer is affirmative.  The source statement and its immediate context
are reproduced below.

\begin{center}
\includegraphics[width=\textwidth]{figures/open_problem_crop.png}
\end{center}

A bounded novelty search on 23 July 2026 used the arXiv id and title, the
exact wording of Open Problem~3.4, and combinations of
\emph{positive contraction}, \emph{monotonically complete},
\emph{Fatou property}, \emph{peripheral point spectrum},
\emph{cyclicity}, and \emph{bilateral shift}.  Searches of the local
2000--2026 arXiv source corpus and web/arXiv-oriented results found the source
paper, its 2016 predecessor \cite{GlueckPeripheral}, and bibliographic mirrors,
but no later work claiming the result below.  This is bounded negative
evidence, not a comprehensive novelty guarantee.

\section*{Proof intuition}

The root-of-unity proof in the source restricts \(T\) to
\(\Fix(T^n)\), where it becomes a lattice automorphism.  For an irrational
unimodular eigenvalue there is no such \(n\).  The replacement is to manufacture
a fixed space on which an auxiliary operator is automatically invertible.

Pass to the bilateral sequence lattice
\(\widehat E=\ell^\infty(\mathbb Z;E)\) and set
\[
 (Sx)_n=T x_{n-1},\qquad (Rx)_n=x_{n+1}.
\]
The operator \(S\) is a positive contraction, so its fixed space \(F\) is a
Banach lattice by the source paper's main theorem.  The bilateral shift \(R\)
commutes with \(S\); hence \(R|_F\) is a lattice automorphism.  An eigenvector
\(Tz=\lambda z\) is encoded by the fixed sequence
\((\lambda^n z)_{n\in\mathbb Z}\), which is a \(\lambda\)-eigenvector of
\(R|_F\).  Cyclicity for lattice automorphisms produces the \(\lambda^k\)
eigenspaces, and evaluation at coordinate zero transfers them back to \(T\).

\section*{Two elementary lemmas}

Recall that a Banach lattice \(E\) has the Fatou property if
\(0\leq x_i\uparrow x\) implies
\(\|x\|=\sup_i\|x_i\|\), and it is monotonically complete if every increasing
norm-bounded net in \(E_+\) has a supremum.

\begin{lemma}[stability under bounded products]\label{lem:product}
If \(E\) is monotonically complete and has the Fatou property, then
\(\ell^\infty(I;E)\), with its coordinatewise order and supremum norm, has
both properties for every index set \(I\).
\end{lemma}

\begin{proof}
Let \((x^\alpha)\) be an increasing norm-bounded net in
\(\ell^\infty(I;E)_+\), with \(\|x^\alpha\|\leq C\).  For each \(i\in I\),
monotone completeness gives
\[
 x_i:=\sup_\alpha x_i^\alpha\in E.
\]
The Fatou property of \(E\) gives
\(\|x_i\|=\sup_\alpha\|x_i^\alpha\|\leq C\), so
\(x=(x_i)_{i\in I}\) belongs to \(\ell^\infty(I;E)\) and is plainly the
supremum of the net.  Finally,
\[
 \|x\|
 =\sup_{i\in I}\sup_\alpha\|x_i^\alpha\|
 =\sup_\alpha\sup_{i\in I}\|x_i^\alpha\|
 =\sup_\alpha\|x^\alpha\|,
\]
which is the Fatou property.
\end{proof}

\begin{lemma}[cyclicity for a lattice automorphism]\label{lem:lattice-auto}
Let \(V\) be a complex Banach lattice and let \(U\) be a lattice
automorphism of \(V\).  For every \(\mu\in\Torus\) and \(k\in\mathbb Z\),
\[
 \dim\ker(\mu-U)\leq\dim\ker(\mu^k-U).
\]
\end{lemma}

\begin{proof}
This is the standard lattice-power argument used in
\cite[Proposition~3.1]{GlueckPeripheral} and invoked by the source in its
proof of Corollary~3.3.  We recall the dimension mechanism.

For \(0\neq f\in V\), represent the principal ideal generated by \(|f|\) as
a \(C(K)\)-space and define the lattice power \(f^{[k]}\) pointwise by
\[
 f^{[k]}(t)=
 \begin{cases}
   \bigl(f(t)/|f(t)|\bigr)^k |f(t)|,&f(t)\neq0,\\
   0,&f(t)=0.
 \end{cases}
\]
This definition is intrinsic.  Since lattice homomorphisms preserve lattice
powers,
\[
 f\in\ker(\mu-U)\setminus\{0\}
 \quad\Longrightarrow\quad
 f^{[k]}\in\ker(\mu^k-U).
\]

For completeness, this implication yields the dimension inequality even
though \(f\mapsto f^{[k]}\) is nonlinear.  Given \(m\) linearly independent
vectors in \(\ker(\mu-U)\), work in the principal ideal generated by the sum
of their moduli.  In its \(C(K)\)-representation, choose \(m\) points and a
new basis \(f_1,\ldots,f_m\) of their span with
\(f_j(t_\ell)=\delta_{j\ell}\).  Then
\(f_j^{[k]}(t_\ell)=\delta_{j\ell}\), including for \(k\leq0\) under the
zero convention above.  Hence the \(f_j^{[k]}\) are linearly independent in
\(\ker(\mu^k-U)\).
\end{proof}

\section*{Full solution of Open Problem 3.4}

\begin{theorem}\label{thm:main}
Let \(E\) be a complex Banach lattice which is monotonically complete and
has the Fatou property, and let \(T\in\cL(E)\) be positive with
\(\|T\|\leq1\).  If \(\lambda\in\Torus\), then for every \(k\in\mathbb Z\),
\[
 \dim\ker(\lambda-T)\leq\dim\ker(\lambda^k-T).
\]
In particular, if \(\lambda\) is an eigenvalue of \(T\), then every integer
power \(\lambda^k\) is an eigenvalue of \(T\).
\end{theorem}

\begin{proof}
Set
\[
 \widehat E:=\ell^\infty(\mathbb Z;E).
\]
By Lemma~\ref{lem:product}, this complex Banach lattice is monotonically
complete and has the Fatou property.  Define \(S,R\in\cL(\widehat E)\) by
\[
 (Sx)_n:=T x_{n-1},
 \qquad
 (Rx)_n:=x_{n+1}
 \qquad(n\in\mathbb Z).
\]
The operator \(S\) is a positive contraction, \(R\) is a lattice isometric
automorphism, and \(SR=RS\).  Let
\[
 F:=\Fix S.
\]
Theorem~3.1 of the source \cite{GlueckFixed} shows that \(F\), with the
inherited order and norm, is a monotonically complete complex Banach lattice
with the Fatou property.  Since both \(R\) and \(R^{-1}\) commute with \(S\),
they leave \(F\) invariant.  Consequently \(R|_F\) and its inverse are
positive, so \(R|_F\) is a lattice automorphism of \(F\).

For each \(\mu\in\Torus\), define
\[
 J_\mu:\ker(\mu-T)\longrightarrow F,\qquad
 (J_\mu z)_n:=\mu^n z.
\]
If \(Tz=\mu z\), then
\[
 (S J_\mu z)_n
 =T(\mu^{n-1}z)
 =\mu^n z,
 \qquad
 (R J_\mu z)_n=\mu^{n+1}z
 =\mu(J_\mu z)_n.
\]
Thus \(J_\mu\) maps into \(\ker(\mu-R|_F)\).  Conversely, if
\(x\in F\) and \(Rx=\mu x\), then \(x_n=\mu^n x_0\), while \(Sx=x\)
at coordinate \(n=1\) gives
\[
 T x_0=x_1=\mu x_0.
\]
Hence evaluation at coordinate zero is the inverse of \(J_\mu\), and
\[
 \dim\ker(\mu-T)=\dim\ker(\mu-R|_F).
\]
Apply Lemma~\ref{lem:lattice-auto} to \(R|_F\):
\[
\begin{aligned}
 \dim\ker(\lambda-T)
 &=\dim\ker(\lambda-R|_F)\\
 &\leq\dim\ker(\lambda^k-R|_F)
 =\dim\ker(\lambda^k-T).
\end{aligned}
\]
This proves the assertion for every \(k\in\mathbb Z\).
\end{proof}

\section*{Additional consequence: the semigroup dimension estimate}

The same dilation also proves the dimension estimate anticipated immediately
after Corollary~3.5 of the source.

\begin{corollary}\label{cor:semigroup}
Let \(E\) be as in Theorem~\ref{thm:main}, and let
\((T_t)_{t\geq0}\) be a positive contractive \(C_0\)-semigroup on \(E\)
with generator \(A\).  Then, for \(\beta\in\mathbb R\) and
\(k\in\mathbb Z\),
\[
 \dim\ker(i\beta-A)\leq\dim\ker(ik\beta-A).
\]
\end{corollary}

\begin{proof}
On \(\widehat E=\ell^\infty(\mathbb R;E)\), define the commuting family of
positive contractions
\[
 (S_t f)(r):=T_t f(r-t)\qquad(t\geq0,\ r\in\mathbb R),
\]
and let \(F=\bigcap_{t\geq0}\Fix S_t\).  Lemma~\ref{lem:product} and
Theorem~3.1 of \cite{GlueckFixed}, now applied to the whole commuting family,
make \(F\) a complex Banach lattice.  Every translation
\((R_u f)(r)=f(r+u)\), \(u\in\mathbb R\), restricts to a lattice
automorphism of \(F\).

The map
\[
 z\longmapsto \bigl(r\mapsto e^{i\beta r}z\bigr)
\]
is a linear bijection from \(\ker(i\beta-A)\) onto
\[
 G_\beta
 :=\{f\in F:R_u f=e^{i\beta u}f\text{ for every }u\in\mathbb R\}.
\]
Indeed, the forward implication follows from
\(T_tz=e^{i\beta t}z\).  Conversely, a member of \(G_\beta\) has the form
\(f(r)=e^{i\beta r}f(0)\), and \(S_t f=f\) at \(r=t\) yields
\(T_t f(0)=e^{i\beta t}f(0)\), which is equivalent to
\(f(0)\in\ker(i\beta-A)\).

If \(0\neq f\in G_\beta\), then every \(R_u\) is a lattice automorphism, so
\[
 R_u(f^{[k]})
 =(R_u f)^{[k]}
 =e^{ik\beta u}f^{[k]}.
\]
Thus \(f^{[k]}\in G_{k\beta}\).  The finite-dimensional separation argument
in Lemma~\ref{lem:lattice-auto}, applied to \(G_\beta\) and \(G_{k\beta}\),
gives \(\dim G_\beta\leq\dim G_{k\beta}\), which is the claimed estimate.
\end{proof}

\section*{Verification notes and limitations}

\begin{itemize}[leftmargin=*]
\item The bilateral construction uses all bounded \(E\)-valued functions.
      There is no measurability, separability, or strong-continuity assumption
      hidden in Lemma~\ref{lem:product}.
\item Both shifts are genuinely bilateral.  This is what makes \(R|_F\)
      invertible with a positive inverse.
\item The sign convention was checked in both directions:
      \(S(\mu^n z)_n=(\mu^n z)_n\), and \(Sx=x\), \(Rx=\mu x\) imply
      \(Tx_0=x_1=\mu x_0\).
\item The only non-elementary input is Theorem~3.1 of the source, together
      with the standard lattice-power dimension lemma already used by the
      source in Corollary~3.3.
\item The result does not extend merely by replacing contractivity with
      power boundedness: the source gives a counterexample under that weaker
      hypothesis.
\item The bounded literature search cannot exclude an unindexed or
      differently phrased proof.  Novelty confidence is moderate.
\end{itemize}

\begin{thebibliography}{9}
\bibitem{GlueckFixed}
J. Gl\"uck,
\emph{A note on the fixed space of positive contractions},
arXiv:2108.05174, 2021.

\bibitem{GlueckPeripheral}
J. Gl\"uck,
\emph{On the peripheral spectrum of positive operators},
Positivity \textbf{20} (2016), 307--336;
arXiv:1503.02430.
\end{thebibliography}

\end{document}
